%%=============================================================================
%% CAPITULO 3 -- PROCEDIMENTOS METODOLOGICOS
%%=============================================================================

\section{Procedimentos metodológicos}

Esta pesquisa adota uma abordagem quantitativa e aplicada, 
utilizando dados do Sistema Nacional de Pesquisa de Custos e Índices da Construção Civil 
(SINAPI) para investigar o potencial de técnicas de ciência de dados e aprendizado 
de máquina na análise e previsão de custos da construção civil. O estudo será desenvolvido 
a partir da organização, tratamento e análise das séries históricas disponíveis na base 
SINAPI, buscando identificar padrões temporais e territoriais que possam contribuir para o 
planejamento de investimentos e para a tomada de decisão em políticas públicas de 
infraestrutura. 

A estratégia metodológica compreenderá etapas de preparação dos dados, 
análise exploratória, construção de variáveis, aplicação de modelos estatísticos e de 
aprendizado de máquina e avaliação dos resultados obtidos. Diferentes técnicas poderão ser 
comparadas por meio de métricas de desempenho adequadas à previsão, permitindo identificar 
os modelos mais aderentes às características dos dados do SINAPI. O algoritmo e o recorte 
temporal definitivos serão estabelecidos durante o desenvolvimento da pesquisa, a partir da 
análise da disponibilidade, qualidade e comportamento das séries históricas utilizadas.

\subsection{Tipo e abordagem da pesquisa}

A pesquisa possui natureza aplicada e abordagem predominantemente quantitativa, 
pois busca produzir conhecimento capaz de apoiar problemas concretos relacionados ao 
planejamento e à gestão de investimentos públicos. A investigação utiliza dados secundários 
provenientes do SINAPI e combina análise exploratória, métodos estatísticos e técnicas de 
aprendizado de máquina para avaliar padrões e possibilidades de previsão dos custos da 
construção civil.

\subsection{Instrumentos de coleta de dados}

A principal fonte de dados será a base do SINAPI, 
considerada no estudo em suas dimensões temporal e territorial.
 Os dados serão preparados para permitir a identificação de inconsistências, valores ausentes, diferenças entre localidades e evolução dos preços e custos ao longo do tempo, constituindo a base analítica para os experimentos realizados na pesquisa.

\subsection{Procedimentos de análise}

A análise será estruturada por meio da comparação entre modelos e seus resultados, 
utilizando métricas quantitativas compatíveis com problemas de previsão. 
A seleção definitiva dos algoritmos será realizada após a etapa exploratória e os experimentos iniciais, 
evitando estabelecer antecipadamente um único método antes de conhecer o comportamento 
dos dados e permitindo avaliar qual abordagem apresenta melhor desempenho para o 
problema investigado.
